% Capitulo 5: Conclusiones (Fase 8, docs/rio_search_plan.md S3.11)
\chapter{Conclusiones}
\label{chap:conclusiones}

\section{Estado al cierre de este cap\'itulo}

Este documento se compila a partir de un banco de pruebas (\emph{Rio\_Search}) ya operativo:
backend con dominio expl\'icito, interfaz web para seguir y comparar corridas, inferencia diaria
del modelo campe\'on y biblioteca de investigaci\'on con exportaci\'on autom\'atica a BibTeX. El
resultado preliminar m\'as relevante -- el BiLSTM supera a la persistencia (\emph{skill}$>0$) en
7 de los 8 horizontes evaluados, con \texttt{t+1} como excepci\'on documentada -- responde
parcialmente a la pregunta de investigaci\'on del Cap\'itulo~\ref{chap:marco-teorico}: es posible
proyectar el caudal de la cuenca alta con un modelo m\'as \'util que la persistencia a partir de
\texttt{t+2}, no a\'un en \texttt{t+1}.

\section{L\'ineas de trabajo futuras}

Quedan como trabajo pendiente para completar esta tesis: la incorporaci\'on del grupo de
\emph{features} de pron\'ostico (ECMWF/GEFS) una vez que llegue a la capa Gold, para responder si
el pron\'ostico meteorol\'ogico mejora la proyecci\'on de caudal m\'as all\'a de lo que ya aporta
el estado hidrol\'ogico observado; la comparaci\'on contra un segundo modelo no basado en redes
profundas (p. ej. \emph{gradient boosting} sobre variables rezagadas); y la promoci\'on a la capa
Gold de las transformaciones experimentales que demuestren aportar valor, siguiendo el mismo
camino de notebook + Decisi\'on documentada que us\'o el resto del pipeline.
